% Source: https://artofproblemsolving.com/community/c3h2283507_foogle_hoogle_bacc_at_it
% Stable identity: aops-topic-2283507-q01
For each prime $p$, let $e_p$ and $f_p$ be the exponents of $p$ in $\gcd(a,b)$ and $\operatorname{lcm}(a,b)$. The exponents of $p$ in $a$ and $b$ are then $\{e_p,f_p\}$ in some order: one of them must be the minimum $e_p$ and the other the maximum $f_p$. If $e_p=f_p$ there is one choice; if $e_p<f_p$ there are two. Choices for different primes are independent.

The primes dividing $21!$ are $2,3,5,7,11,13,17,19$, and $21=3\cdot7$. For each of these $8$ primes the exponent in $21!$ is larger than in $21$: for $3$ and $7$ the exponents in $21!$ are $7+2=9>1$ and $3>1$, and the other six primes do not divide $21$ at all. So every one of the $8$ primes gives two choices, and the number of ordered pairs is $2^8=\boxed{256}$.
